\section{Claim-to-proof and evidence map}
\label{hfmap:main}

Every mathematical assertion in the analysis is either a definition, an
explicit model hypothesis, or a consequence covered below. A calculated
decimal is an approximation to a proved exact expression; a citation to an
empirical finding or a recorded software run is evidence of that kind,
not a mathematical theorem about people or arbitrary programs. This
distinction is necessary to make the scope of ``all claims'' precise.

\begingroup\small
\setlength{\tabcolsep}{4pt}
\begin{longtable}{>{\raggedright\arraybackslash}p{.35\linewidth}>{\raggedright\arraybackslash}p{.59\linewidth}}
\toprule
Claim/location & Proof, hypothesis or evidence and limitation\\
\midrule
\endfirsthead
\toprule Claim/location & Proof, hypothesis or evidence and limitation\\\midrule
\endhead
\bottomrule\endfoot
Original article's variance identity, equation~\ref{eq:variance-sum};
finite variances suffice &
Proposition~\ref{hfpr:standardization}, by covariance expansion and
Cauchy--Schwarz; no normality is assumed. Attribution of the identity and
the selected 27 scorers is to \citet{hoeflin}.\\
Standardization and the SD-15/SD-16 examples,
equation~\ref{eq:standardize} &
Proposition~\ref{hfpr:standardization}. Numerical \(z=4\) and
\(z=15/4\) are exact rational calculations. Population means, SDs and
comparability of norms must be supplied as assumptions/evidence.\\
Weighted variance, composite and its mean/variance,
equations~\ref{eq:weighted-variance}--\ref{eq:weighted-composite} &
Proposition~\ref{hfpr:standardization}; weights are fixed and
\(w^{\mathsf T}Rw>0\). If estimated parameters replace population
values, the asserted population identity no longer follows automatically;
the equations themselves identify the unknown inputs.\\
Average correlation, mean variance, equal-weight formula and singleton,
equations~\ref{eq:average-correlation}--\ref{eq:equal-composite} &
Proposition~\ref{hfpr:equal-weights}. This includes the PSD domain,
the zero-denominator boundary, and a constructive demonstration that the
full matrix matters to other weights.\\
Affine reporting units and common-scale IQ formula,
equations~\ref{eq:reporting-scale}--\ref{eq:iq-composite} &
Propositions~\ref{hfpr:standardization} and \ref{hfpr:equal-weights}.
These are scale definitions with proved population moments, rather than
an empirical validation of an IQ interpretation.\\
Three-score worked example, equation~\ref{eq:worked-example} &
Proposition~\ref{hfpr:numerics} gives exact radicals and rational
squared bounds proving its displayed rounding.\\
Moments and marginal normality do not establish a Gaussian composite;
joint Gaussianity suffices &
Proposition~\ref{hfpr:normality}, including an explicit normal-marginal,
same-correlation counterexample with an atom in the sum. The general
Gaussian foundation is Proposition~\ref{pf:gaussian}.\\
Percentile, normal quantile and required average cutoff,
equations~\ref{eq:quantile}--\ref{eq:threshold} &
Proposition~\ref{hfpr:thresholds} proves the quantile definition and
event equivalence. The decimal quantile is a computed approximation
documented by \citet{rnormal} and the calculation script; it is not an
assumption-free empirical frequency.\\
All rows of Table~\ref{tab:thresholds}, table corrections, fixed composite
cutoff, SD-15 singleton and correlation-0.90 example &
Proposition~\ref{hfpr:numerics} supplies each exact expression and its
quantile-error propagation. \texttt{data/thresholds.csv} and
\texttt{scripts/verify\_calculations.py} reproduce the decimals.
Quantile-dependent last-digit rounding is numerical evidence, not a
certified real-arithmetic enclosure. Operational rounding is a scoring
policy, not a probability theorem.\\
Limit, absence of a score ceiling, effects of changing correlation,
equation~\ref{eq:limit} &
Proposition~\ref{hfpr:thresholds}, including derivative signs, an error
bound, admissibility of the fixed-correlation sequence, and unbounded
Gaussian support. No claim is made for arbitrary changing batteries.\\
Selection may change correlations and no correction is universal &
Proposition~\ref{hfpr:selection} gives a fully derived half-normal
selection example. Actual range-restriction corrections are attributed
to \citet{sackett2000}; no empirical mechanism is inferred from the
counterexample.\\
Selected small samples do not establish a one-in-a-million frequency;
repeated/best-score/stopping procedures need calibration &
Proposition~\ref{hfpr:selection} proves finite-sample nonidentification
and independent-attempt versus perfectly dependent counterexamples.
``Can differ'' is the precise claim; perfectly dependent procedures
can coincide. Practice effects are the empirical finding of
\citet{hausknecht2007}.\\
Population score variation differs from individual uncertainty;
\(R\) alone does not identify a latent target &
Propositions~\ref{hfpr:factor} and \ref{hfpr:nonidentification};
two models have identical observable scores and different latent
posteriors. Measurement validity and extreme-score precision require
additional evidence, not another variance identity.\\
One-factor assumptions and all latent formulas,
equations~\ref{eq:factor-model}--\ref{eq:latent-variance} &
The model is expressly hypothetical. Proposition~\ref{hfpr:factor}
proves its covariance, reliability-like squared correlation, posterior
normal law, contrasts' lack of extra information and conditional-mean
MSE optimality. Proposition~\ref{hfpr:numerics} gives exact expressions
for all accompanying numerical approximations.\\
Recommendations about battery design, norms, validity and transparent
reporting &
Normative proposals, explicitly attributed to the testing standards
\citep[Standards 1.1--1.4, 1.14, 5.8--5.9]{standards2014}.
A normative recommendation is not presented as a theorem derived from
the common-factor model.\\
Numerical reproduction and the software-verification appendix &
Reproducible numerical evidence and source-pinned historical campaign
records, as linked in Section~\ref{sec:verification}. Those records,
rather than a mathematical citation, support check totals, source
coverage, negative controls, replay, tool versions and contract counts.
The historical counts describe their snapshots, not the current tree.\\
What software evidence does and does not establish &
Native SPARK checks concern the encoded program contracts under trusted
interfaces; Deal runtime/property tests, limited symbolic checks,
structural counts, numerical comparisons and model validity remain
separate. For the mathematical sampling claims see
Proposition~\ref{hfpr:simulation} and
Section~\ref{sec:simulation-statistical-methods}. No proof of foreign
libraries, pseudorandom independence, physical I/O or human-score
Gaussianity is claimed.\\
AI contribution, provenance and absence of participant data &
Project process disclosure in Section~\ref{sec:ai-disclosure}, not a
mathematical proposition or a claim of independent human review.\\
\end{longtable}
\endgroup
